\chapter{Considerações Finais}
\label{cap:consideracao}

Diante do exposto, este Trabalho de Conclusão de Curso 1 fundamentou teoricamente, estruturou e validou a viabilidade da arquitetura neuro-simbólica proposta. Ao longo desta etapa, foi delineada uma esteira automatizada de triagem de alertas de segurança e estabelecido um plano metodológico experimental para avaliar o impacto do enriquecimento semântico na filtragem de falsos positivos no ecossistema da linguagem Go. A prova de conceito desenvolvida demonstrou a viabilidade técnica do ecossistema de simulação e do \textit{middleware} de hidratação de contexto, demonstrando a execução do fluxo de ponta a ponta e evidenciando a relevância prática de separar a cobertura do motor simbólico do acerto da camada neural. 

Para a continuidade da pesquisa no Trabalho de Conclusão de Curso 2, o foco central consistirá na execução do experimento em larga escala. As atividades subsequentes englobam a expansão da base de validação para a totalidade das amostras do \textit{dataset} SastBench e das CVEs catalogadas na OSV.dev, além da execução do desenho fatorial completo. Essa próxima fase contrastará os modelos de fronteira selecionados sob as abordagens de \textit{prompts} de controle (\textit{baseline}) e especialista, permitindo a extração das métricas estatísticas propostas --- como o Coeficiente de Correlação de Matthews (MCC) e a Taxa de Redução de Alertas (TRA) --- e a auditoria qualitativa das justificativas estruturadas da IA. Por fim, os desafios operacionais identificados no piloto, como a gestão de cotas de API, serão mitigados por meio de estratégias de processamento incremental em lotes, culminando na validação empírica da solução e na resposta definitiva às questões de pesquisa formuladas.